\section{Líneas futuras}

De todo lo que el trabajo deja abierto, tres líneas destacan sobre las demás por ser las que
condicionan que esta plataforma pueda llegar a volar.

\textbf{Completar SpaceWire, desde la capa física hasta un controlador propio.} Es la única
interfaz de las placas que no llegó a validarse. La tarjeta AOCS expone los dos enlaces con su
señalización LVDS, sus repetidores y sus pares igualados en longitud, de modo que el trabajo
pendiente está íntegramente en la lógica programable y el software: la codificación
\textit{data-strobe}, la máquina de estados de establecimiento y recuperación del enlace, y los
niveles de paquete y de red del estándar. El camino razonable es el mismo que siguió el transceptor
serie de este trabajo: implementar el controlador desde cero en VHDL, verificarlo por simulación
bloque a bloque y exponerlo al PS mediante el transporte MCDMA ya desarrollado, que admite canales
adicionales sin rediseño. Es el paso más relevante a corto plazo, porque SpaceWire es el protocolo
previsto para los enlaces de mayor ancho de banda del satélite.

\textbf{Endurecer la variante MCDMA mediante triplicación modular (TMR).} La lógica programable
comercial es susceptible a los efectos de la radiación en órbita, en particular a los
\textit{single event upsets} (SEU), que alteran el contenido de los biestables y de la memoria de
configuración. La técnica habitual de mitigación es la triplicación modular redundante: tres copias
de la lógica y módulos votadores que enmascaran el fallo de cualquiera de ellas. La línea de
trabajo concreta es replicar el proyecto de la variante MCDMA triplicando su lógica propia (el
puente VHDL y los transceptores serie), incorporar los votadores necesarios y demostrar que el
sistema sigue superando la misma campaña de treinta y seis comprobaciones. El presupuesto de área
lo permite, y esa fue una de las razones de escoger esta arquitectura: triplicar el 5,2\,\% de
ocupación del MCDMA mantiene el diseño en torno al 16\,\% del dispositivo, mientras que la misma
operación sobre la variante DMA14 rondaría el 44\,\% y la haría inviable
(sección~\ref{sec:resultados_benchmark}).

\textbf{Proteger la memoria de configuración de la lógica programable.} La triplicación cubre la
lógica de usuario y deja fuera un punto que en una FPGA basada en SRAM es igual de sensible: la
memoria de configuración. Un SEU en ella no corrompe un dato, sino el propio circuito
implementado, y la triplicación no lo detecta porque las tres copias residen en el mismo
dispositivo. El mecanismo estándar frente a esto es el \textit{scrubbing}, la reescritura
periódica de la memoria de configuración a partir de una copia de referencia, que en el Zynq
UltraScale+ puede gobernarse desde el PS. La línea de trabajo consiste en integrarlo sobre la
variante MCDMA y medir dos cosas: cada cuánto puede refrescarse la configuración sin perturbar el
tráfico de los canales, y qué parte del área queda cubierta.

Fuera ya del alcance de una continuación inmediata queda la comparación sistemática de mecanismos
de tolerancia a radiación sobre esta plataforma, que constituye un trabajo en sí mismo. La
migración a memorias externas con código corrector o la evaluación de dispositivos tolerantes a
radiación por construcción pertenecen a esa categoría y condicionan el diseño de la placa entera,
no solo el de la lógica desarrollada aquí.
